\documentclass[conference]{IEEEtran}
\IEEEoverridecommandlockouts

\usepackage[utf8]{inputenc}
\usepackage[T1]{fontenc}
\usepackage[spanish]{babel}
\usepackage{cite}
\usepackage{amsmath,amssymb}
\usepackage{graphicx}
\usepackage{url}
\usepackage{hyperref}
\usepackage{booktabs}
\usepackage{listings}
\usepackage{xcolor}

\lstset{
  basicstyle=\ttfamily\footnotesize,
  breaklines=true,
  frame=single,
  language=C,
  keywordstyle=\color{blue},
  commentstyle=\color{gray},
}

\begin{document}

\title{perfanalyzer: Una API en C para el Análisis de Rendimiento\\
de Contenedores Docker}

\author{
  \IEEEauthorblockN{Camilo Loaiza Hernández}
  \IEEEauthorblockA{Facultad de Ingeniería\\
  Universidad de Antioquia\\
  Medellín, Colombia\\
  camiloloaizah03@gmail.com}
  \and
  \IEEEauthorblockN{[Nombre compañero]}
  \IEEEauthorblockA{Facultad de Ingeniería\\
  Universidad de Antioquia\\
  Medellín, Colombia\\
  correo@udea.edu.co}
}

\maketitle

% ─── Abstract ────────────────────────────────────────────────────────────────
\begin{abstract}
Este artículo presenta \textit{perfanalyzer}, una API implementada en C que
expone métricas de CPU, memoria y E/S de procesos tanto desde dentro de un
contenedor Docker (vía \texttt{/proc}) como desde el host (vía cgroup v2). La
API incluye un servidor HTTP/REST y un sistema de \textit{profiling} por función
instrumentado con macros. Se realizó un diseño experimental con $n=30$ réplicas
por celda evaluando tres tipos de carga de trabajo (CPU, memoria, E/S) bajo
tres configuraciones de contenedor y modo host, con análisis estadístico
(Welch t-test / Mann-Whitney, IC 95\%, Cohen's d). Los resultados muestran
[COMPLETAR CON RESULTADOS REALES].
\end{abstract}

\begin{IEEEkeywords}
Docker, cgroup v2, rendimiento, profiling, API REST, C, análisis estadístico
\end{IEEEkeywords}

% ─── I. Introducción ─────────────────────────────────────────────────────────
\section{Introducción}

Los contenedores han transformado el despliegue de software al ofrecer
aislamiento ligero y reproducibilidad. Sin embargo, el análisis de su
rendimiento presenta desafíos específicos: las herramientas tradicionales
como \texttt{perf} o \texttt{gprof} no tienen visibilidad directa de los
límites impuestos por el scheduler y la jerarquía de control de grupos
(\textit{cgroups}) \cite{sobieraj2024docker}.

La motivación de este proyecto es construir una herramienta de propósito
educativo e investigativo que: (a) sea implementada en lenguaje C sobre Linux,
sin dependencias de alto nivel; (b) use exclusivamente las interfaces del kernel
(\texttt{/proc}, cgroup v2) para obtener métricas; (c) las exponga a través de
una API REST; y (d) permita \textit{profiling} por función con sobrecarga mínima.

La metodología consiste en diseñar la librería \textit{perfanalyzer},
instrumentar tres cargas de trabajo representativas y comparar su rendimiento
en host versus tres configuraciones de Docker, validando estadísticamente las
diferencias observadas.

% ─── II. Marco Teórico ───────────────────────────────────────────────────────
\section{Marco Teórico}

\subsection{cgroup v2 y aislamiento de recursos}

El subsistema \textit{control groups} (cgroups) v2, estabilizado en Linux 4.5,
organiza los procesos en una jerarquía unificada bajo
\texttt{/sys/fs/cgroup}. Cada cgroup expone archivos como
\texttt{cpu.stat} (tiempo de CPU), \texttt{memory.current} y
\texttt{memory.peak} (uso de memoria) e \texttt{io.stat} (bytes por
dispositivo) \cite{docker2022monitor}.

Docker utiliza cgroups para implementar los límites de \texttt{--cpus} y
\texttt{--memory}. Un contenedor con \texttt{--cpus=1} tiene su tiempo de CPU
restringido por el parámetro \texttt{cpu.max} del cgroup, lo que reduce el
\textit{throughput} de cargas CPU-intensivas.

\subsection{Profiling por función en espacio de usuario}

El \textit{profiling} de alta resolución en espacio de usuario se realiza
con \texttt{clock\_gettime(CLOCK\_MONO\-TONIC\_RAW)}, que es inmune al ajuste
NTP \cite{intel2024profiling}. La combinación con \texttt{getrusage(RUSAGE\_THREAD)}
permite aislar el tiempo de CPU consumido específicamente por el hilo del
\textit{probe}, separando el tiempo de E/S y espera.

\subsection{Métodos estadísticos}

Para comparar dos distribuciones se aplica primero la prueba de normalidad
Shapiro-Wilk (adecuada para $n \leq 5000$) \cite{ieee2024performance}. Si ambas
distribuciones son normales se usa la t de Welch (no asume varianzas iguales);
en caso contrario, Mann-Whitney U (no paramétrica). El tamaño del efecto se
reporta como Cohen's $d$ o $r$ de rango biserial. Los intervalos de confianza
al 95\% se calculan por bootstrap ($n = 2000$ remuestras).

% ─── III. Metodología ────────────────────────────────────────────────────────
\section{Metodología}

El proyecto siguió un proceso iterativo de cuatro etapas: (1) diseño de la
arquitectura y la API pública; (2) implementación de los módulos de muestreo
y profiling; (3) implementación del servidor HTTP/REST; (4) diseño y ejecución
del experimento con análisis estadístico.

El entorno de desarrollo fue Windows 11 con VSCode y Claude Code como
asistente; el entorno de ejecución y pruebas fue una VM Ubuntu 24.04 LTS
con cgroup v2 nativo, Docker 24+ y GCC 11.

% ─── IV. Implementación ──────────────────────────────────────────────────────
\section{Implementación}

\subsection{Arquitectura general}

\textit{perfanalyzer} se estructura en tres componentes (Figura~\ref{fig:arch}):

\begin{enumerate}
  \item \textbf{libperfanalyzer}: librería estática con los módulos de muestreo
    (\texttt{sampler\_cpu}, \texttt{sampler\_mem}, \texttt{sampler\_io},
    \texttt{sampler\_cgroup}), el profiler (\texttt{profiler.c}) y el exportador
    JSON (\texttt{exporter\_json.c}).
  \item \textbf{perfanalyzer\_server}: daemon HTTP/REST que expone ocho
    endpoints sobre la librería.
  \item \textbf{perfctl}: cliente CLI para interactuar con el servidor o tomar
    muestras directas.
\end{enumerate}

\begin{figure}[h]
  \centering
  % \includegraphics[width=\columnwidth]{arch.pdf}
  \fbox{\parbox{0.9\columnwidth}{\centering
    [Figura de arquitectura — agregar diagrama]}}
  \caption{Arquitectura de perfanalyzer.}
  \label{fig:arch}
\end{figure}

\subsection{Módulos de muestreo}

\textbf{CPU}: Se leen \texttt{/proc/self/stat} y \texttt{/proc/stat} en dos
instantes sucesivos. El porcentaje de CPU del proceso se calcula como:
\begin{equation}
  \%\mathrm{CPU} = \frac{\Delta(\mathit{utime} + \mathit{stime}) / \mathit{ticks\_per\_s}}
                        {\Delta t_{\mathrm{wall}}} \times 100
\end{equation}

\textbf{Memoria}: Se parsea \texttt{/proc/self/status} buscando
\texttt{VmRSS}, \texttt{VmSize} y \texttt{VmPeak}.

\textbf{E/S}: Se lee \texttt{/proc/self/io} para obtener
\texttt{read\_bytes}, \texttt{write\_bytes}, \texttt{syscr} y \texttt{syscw}.

\textbf{cgroup v2}: El path propio se detecta leyendo \texttt{/proc/self/cgroup}
(línea \texttt{0::}). Luego se leen \texttt{cpu.stat}, \texttt{memory.current},
\texttt{memory.peak} e \texttt{io.stat}.

\subsection{Profiler por función}

La macro \texttt{PERF\_PROBE(ctx, "nombre")} usa
\texttt{\_\_attribute\_\_((cleanup))} para garantizar que \texttt{perf\_probe\_end}
se invoque aun ante retornos anticipados:

\begin{lstlisting}
PERF_PROBE(ctx, "matmul") {
    matmul(C, A, B, n);
}
\end{lstlisting}

Los percentiles se estiman por \textit{reservoir sampling} (Algorithm R de
Vitter, $k = 1024$) para mantener uso de memoria constante.

\subsection{Servidor HTTP/REST}

El servidor implementa un parser HTTP/1.1 mínimo con un hilo por conexión
(detached). Los endpoints principales son:

\begin{itemize}
  \item \texttt{GET /snapshot} — muestra puntual.
  \item \texttt{POST /recording/\{start,stop\}} — serie temporal.
  \item \texttt{GET /probes} — estadísticas de sondas.
\end{itemize}

La respuesta es siempre JSON plano generado sin dependencias externas.

\subsection{Cargas de trabajo}

Se implementaron tres cargas instrumentadas con \texttt{PERF\_PROBE}:
\textbf{cpu\_stress} (multiplicación de matrices $N \times N$ + SHA-256),
\textbf{mem\_stress} (alloc/touch/free + stride walk) e
\textbf{io\_stress} (escritura y lectura secuencial y aleatoria de archivos).
Cada una acepta \texttt{--intensity 1|2} y \texttt{--duration S}.

% ─── V. Protocolo de Experimentación ─────────────────────────────────────────
\section{Protocolo de Experimentación}

\subsection{Variables}

\textbf{Independientes}: entorno (\textit{host}, \textit{docker:baseline},
\textit{docker:cpu-limit}, \textit{docker:mem-limit}); tipo de carga (CPU,
Mem, IO); intensidad (1=baja, 2=alta).

\textbf{Dependientes}: CPU promedio (\%), RSS promedio (kB), bytes de E/S,
latencia p99 de sonda (µs), \textit{throughput} (iteraciones en 10 s).

\subsection{Diseño}

$n = 30$ réplicas por celda; 2 corridas de calentamiento descartadas;
orden aleatorizado dentro de cada bloque. El experimento se automatizó con
\texttt{run\_experiment.sh}.

\subsection{Validación cruzada}

Para validar la precisión de los samplers se comparó:
\texttt{/snapshot} (interno) vs.\ \texttt{docker stats} (host) vs.\ cgroup
desde el host. La diferencia relativa admisible es $< 5\%$ en CPU promedio.

\subsection{Herramientas}

\texttt{GCC 11} + CMake + ASan/UBSan (correctness), \texttt{valgrind}
(fugas de memoria), Python + NumPy/SciPy/pandas/matplotlib (análisis).
La elección de Python para análisis se justifica por la disponibilidad de
\texttt{scipy.stats} con implementaciones validadas de los tests estadísticos.

% ─── VI. Resultados ───────────────────────────────────────────────────────────
\section{Resultados}

% ⚠️ COMPLETAR con datos reales del experimento.
% Sustituir los valores de ejemplo por los del CSV.

\subsection{Overhead del sampler}

La diferencia de CPU promedio entre la misma carga con y sin sampler activo
fue de $\Delta\mu = \mathbf{X.X\%}$ (IC 95\%: [$a$, $b$]), lo que indica
un overhead [despreciable / moderado].

\subsection{CPU-stress: host vs. Docker}

La Tabla~\ref{tab:cpu} muestra los resultados para \texttt{cpu\_stress}
con intensidad alta. [COMPLETAR]

\begin{table}[h]
  \centering
  \caption{CPU promedio (\%) — cpu\_stress (intensidad alta, n=30)}
  \label{tab:cpu}
  \begin{tabular}{lcccc}
    \toprule
    Entorno & Media & IC 95\% & Test & p-valor \\
    \midrule
    host           & -- & [--,--] & -- & -- \\
    docker:baseline & -- & [--,--] & Welch & -- \\
    docker:cpu-limit & -- & [--,--] & Welch & -- \\
    docker:mem-limit & -- & [--,--] & Welch & -- \\
    \bottomrule
  \end{tabular}
\end{table}

\subsection{Memoria y E/S}

[COMPLETAR con tablas y figuras de mem\_stress e io\_stress.]

\subsection{Figuras}

\begin{figure}[h]
  \centering
  % \includegraphics[width=\columnwidth]{figures/boxplot_cpu_pct_avg.png}
  \fbox{\parbox{0.9\columnwidth}{\centering
    [Boxplot CPU — agregar figura del experimento]}}
  \caption{Distribución de CPU promedio por entorno y carga.}
  \label{fig:boxplot_cpu}
\end{figure}

% ─── VII. Conclusiones ────────────────────────────────────────────────────────
\section{Conclusiones}

[COMPLETAR tras obtener resultados.]

Entre las lecciones generales:
\begin{itemize}
  \item La lectura de \texttt{/proc} introduce latencia de orden $O(\mu s)$ por
    muestra, lo que es admisible para períodos $\geq 10\,\mathrm{ms}$.
  \item cgroup v2 provee métricas más consistentes que \texttt{/proc} bajo
    contenedores con límites de CPU, pues el scheduler los aplica a nivel del
    cgroup.
  \item El \textit{reservoir sampling} con $k=1024$ estima percentiles con
    error $< 5\%$ para $n \geq 30$.
\end{itemize}

% ─── Referencias ─────────────────────────────────────────────────────────────
\begin{thebibliography}{9}

\bibitem{sobieraj2024docker}
M. Sobieraj y D. Kotynski,
``Docker performance evaluation across operating systems,''
\emph{Applied Sciences}, vol. 14, no. 15, p. 6672, 2024.
\url{https://doi.org/10.3390/app14156672}

\bibitem{docker2022monitor}
Docker Inc.,
``How to monitor container memory and CPU usage in Docker Desktop,''
Docker Blog, 2022.
\url{https://www.docker.com/blog/how-to-monitor-container-memory-and-cpu-usage-in-docker-desktop/}

\bibitem{intel2024profiling}
Intel,
``Profiling Docker containers,''
Intel VTune Profiler Cookbook 2024.
\url{https://www.intel.com/content/www/us/en/docs/vtune-profiler/cookbook/2024-1/profiling-in-docker-container.html}

\bibitem{ieee2024performance}
IEEE,
``Performance analysis of container technologies for computer systems,''
\emph{IEEE Xplore}, 2024.
\url{https://ieeexplore.ieee.org/document/10466754/}

\bibitem{ieee2024linux}
IEEE,
``Performance of Docker containerization tool on different Linux distributions,''
\emph{IEEE Xplore}, 2024.
\url{https://ieeexplore.ieee.org/document/10496034/}

\bibitem{ieee2023arm}
IEEE,
``Benchmarking container technologies on ARM-based edge devices,''
\emph{IEEE Xplore}, 2023.
\url{https://ieeexplore.ieee.org/document/10268421/}

\bibitem{arxiv2024vm}
arXiv,
``Experimental assessment of containers running on top of virtual machines,''
2024.
\url{https://arxiv.org/abs/2401.07539}

\bibitem{scidir2024redb}
ScienceDirect,
``REDB: Real-time enhancement of Docker containers via memory,''
\emph{Journal of Systems Architecture}, 2024.
\url{https://www.sciencedirect.com/science/article/abs/pii/S1383762124000729}

\bibitem{scidir2023leakage}
ScienceDirect,
``A framework for evaluating information leakages of Docker containers,''
\emph{SoftwareX}, 2023.
\url{https://www.sciencedirect.com/science/article/pii/S2352711023002728}

\end{thebibliography}

\end{document}
